% !TeX program = pdflatex
\documentclass[11pt,a4paper]{article}

\usepackage[utf8]{inputenc}
\usepackage[T1]{fontenc}
\usepackage[french]{babel}
\usepackage{geometry}
\geometry{margin=2.4cm}

\usepackage{graphicx}
\usepackage{float}
\usepackage{xcolor}
\usepackage{hyperref}
\usepackage{listings}
\usepackage{enumitem}
\usepackage{caption}

\hypersetup{
  colorlinks=true,
  linkcolor=blue,
  urlcolor=blue
}

\setlist[itemize]{topsep=4pt,itemsep=2pt}

\lstset{
  basicstyle=\ttfamily\small,
  breaklines=true,
  frame=single,
  rulecolor=\color{black!15},
  backgroundcolor=\color{black!2}
}

% ====== À REMPLIR PAR VOUS ======
\newcommand{\ProjectTitle}{UniSmartPay — Système de paiement universitaire RFID + code-barres}
\newcommand{\GroupClass}{\textit{(Classe / groupe à compléter)}}
\newcommand{\MemberA}{Ele Meskini}
\newcommand{\MemberB}{Mohamed Khalil Jadlaoui}
\newcommand{\MemberC}{Fayech Jallel}
\newcommand{\MemberD}{Fourat Andolsi}
\newcommand{\MemberE}{Ahmed Smaili}
\newcommand{\GitHubURL}{\url{https://github.com/ORG/REPO}}

% Helper: safe includegraphics (compile even if file missing)
% Usage: \SmartFigure{figures/file.png}{0.9\linewidth}
\newcommand{\SmartFigure}[2]{%
  \IfFileExists{#1}{\includegraphics[width=#2]{#1}}{%
    \fbox{\begin{minipage}[c][0.18\textheight][c]{#2}
      \centering\small\textit{Image manquante}\\
      \texttt{#1}\\
      \vspace{4pt}
      Déposez l'image dans \texttt{report/figures/}
    \end{minipage}}%
  }%
}

\title{\ProjectTitle}
\author{\MemberA \and \MemberB \and \MemberC \and \MemberD \and \MemberE}
\date{Avril 2026}

\begin{document}

\begin{titlepage}
  \centering

  % Logos (optionnels)
  \begin{figure}[H]
    \centering
    \begin{minipage}{0.48\textwidth}
      \centering
      \SmartFigure{figures/logo_fac.png}{0.55\linewidth}
    \end{minipage}
    \begin{minipage}{0.48\textwidth}
      \centering
      \SmartFigure{figures/logofst.png}{0.55\linewidth}
    \end{minipage}
  \end{figure}

  \vspace{8pt}
  {\LARGE\bfseries \ProjectTitle\par}
  \vspace{8pt}
  {\large Document technique du projet (Rendu final)\par}
  \vspace{10pt}

  {\GroupClass\par}
  \vspace{12pt}

  \begin{tabular}{ll}
    \textbf{Groupe :} & \MemberA, \MemberB, \MemberC, \MemberD, \MemberE\\
    \textbf{Date :} & \today\\
    \textbf{Dépôt GitHub :} & \GitHubURL\\
  \end{tabular}

  \vfill
  \begin{minipage}{0.9\textwidth}
    \small
    \textbf{Résumé.} L’idée principale est de rendre la \textbf{carte étudiante “smart”} : elle devient une clé d’accès à un compte prépayé consultable, rechargeable et traçable, utilisable pour payer rapidement à la buvette et au restaurant. UniSmartPay combine la lecture RFID (carte étudiante) et le scan de code-barres (produits), avec une interface web (admin/étudiant/terminal), une base MySQL et un terminal ESP32/RC522.
  \end{minipage}
\end{titlepage}

\tableofcontents
\newpage

\section{Informations groupe et projet}
\begin{itemize}
  \item \textbf{Projet :} UniSmartPay (paiement campus RFID + code-barres).
  \item \textbf{Idée principale :} la carte étudiante devient une carte de paiement intelligente (identification + accès au solde + paiement).
  \item \textbf{Objectif :} réduire les files d’attente et assurer la traçabilité (référence + solde avant/après).
  \item \textbf{Composants :} Web PHP (PDO), MySQL, scan code-barres via webcam, terminal ESP32 + lecteur RFID RC522.
\end{itemize}

\subsection{Répartition des rôles (résumé)}
\begin{itemize}
  \item \textbf{\MemberB} : ESP32/RC522, câblage, modes terminal, intégration API, scan/caméra.
  \item \textbf{\MemberE} : design web / UI-UX.
  \item \textbf{\MemberC} : backend PHP (API + logique métier) + MySQL (transactions).
  \item \textbf{\MemberA} : backend PHP (pages/validations) + rapport/documentation + livrables.
  \item \textbf{\MemberD} : \textbf{pas de tâches PHP} ; aide \MemberB sur le \textbf{code ESP32} (tests/debug/logs) + validation/jeux de données.
\end{itemize}

\section{Lien GitHub}
\GitHubURL

\section{Instructions pour exécuter le projet}
\subsection{Prérequis}
\begin{itemize}
  \item Windows + XAMPP (Apache + MySQL).
  \item PHP 8.x recommandé.
  \item Navigateur Chrome/Edge pour le Web Serial (optionnel) et le scan webcam.
\end{itemize}

\subsection{Installation (XAMPP + MySQL) — étapes détaillées}
\begin{enumerate}
  \item Copier le dossier du projet dans \texttt{htdocs} (ex: \texttt{C:\\xampp\\htdocs\\UniSmartPay}).
  \item Démarrer XAMPP : \textbf{Apache} et \textbf{MySQL} en \textit{Running}.
  \item Ouvrir phpMyAdmin : \url{http://localhost/phpmyadmin}.
  \item Importer la base de données : \texttt{database/unismart_pay_schema.sql}.
  \begin{itemize}
    \item Si le script SQL crée déjà la base, il suffit de \textit{Importer}.
    \item Sinon : créer une base (ex: \texttt{unismart\_pay}) puis importer le fichier SQL dans cette base.
  \end{itemize}
  \item Vérifier / adapter la configuration DB dans \texttt{web/config/config.php} (ex: \texttt{DB\_HOST=127.0.0.1}, \texttt{DB\_PORT=3307}, \texttt{DB\_USER=root}).
\end{enumerate}

\subsection{Configuration (paramètres importants)}
\begin{itemize}
  \item \textbf{MySQL Port :} dans ce projet, MySQL est souvent sur \texttt{3307}. Assurez-vous que \texttt{DB\_PORT} correspond à votre XAMPP.
  \item \textbf{DB\_HOST :} en cas d’erreur de connexion, utiliser \texttt{127.0.0.1} (force TCP) au lieu de \texttt{localhost}.
  \item \textbf{URL d’accès :} l’application web tourne sur \texttt{http://localhost/UniSmartPay/web/}.
\end{itemize}

\subsection{Vérification rapide (smoke test)}
\begin{enumerate}
  \item Ouvrir l’accueil : \url{http://localhost/UniSmartPay/web/}.
  \item Se connecter en admin et vérifier que le dashboard s’affiche.
  \item Vérifier qu’un endpoint JSON répond (exemple) : \texttt{/web/api/card\_check.php}.
\end{enumerate}

\subsection{Lancer l’application}
\begin{itemize}
  \item Accueil : \url{http://localhost/UniSmartPay/web/}
  \item Admin : \url{http://localhost/UniSmartPay/web/admin/login.php}
  \item Étudiant : \url{http://localhost/UniSmartPay/web/etudiant/login.php}
\end{itemize}

\subsection{Comptes de démonstration}
\begin{itemize}
  \item Admin : \texttt{admin@unismart.tn} / \texttt{Admin@1234}
  \item Étudiant : \texttt{2024/FST/0001} / \texttt{Etudiant@1234}
\end{itemize}

\subsection{Données minimales pour une démonstration (checklist)}
Pour valider le parcours complet (scan + paiement), préparer au minimum :
\begin{itemize}
  \item 1 étudiant avec un compte et un \textbf{solde} positif.
  \item 1 carte RFID \textbf{assignée} à cet étudiant et \textbf{active}.
  \item 3 produits avec \textbf{code-barres} et \textbf{prix} (pour le mode BUVETTE).
\end{itemize}
\textbf{Conseil :} créer ces données via l’interface admin (plus simple à justifier dans un rendu que des inserts SQL).

\subsection{Optionnel : terminal RFID (ESP32)}
\begin{itemize}
  \item Firmware : \texttt{esp32/unismart_pay_terminal.ino}
  \item Endpoints principaux :
  \begin{itemize}
    \item Legacy (paiement RESTO) : \texttt{/web/api/payment.php}
    \item Flux buvette (commande) : \texttt{/web/api/terminal\_order\_claim.php} et \texttt{/web/api/terminal\_order\_pay.php}
    \item Vérification carte : \texttt{/web/api/card\_check.php}
    \item Mode dynamique : \texttt{/web/api/terminal\_mode\_get.php} et \texttt{/web/api/terminal\_mode\_set.php}
  \end{itemize}
\end{itemize}

\subsubsection*{Point critique réseau (ESP32 \texorpdfstring{$\leftrightarrow$}{<->} serveur)}
\begin{itemize}
  \item Depuis l’ESP32, \texttt{http://localhost/...} ne fonctionne pas (localhost = ESP32).
  \item Utiliser l’\textbf{IP locale} du PC qui héberge XAMPP (ex: \texttt{http://192.168.1.10/UniSmartPay/web/api/...}).
  \item En cas de timeout : vérifier que le PC et l’ESP32 sont sur le même WiFi et que le firewall autorise Apache.
\end{itemize}

\subsection{Déroulé de démonstration (recommandé)}
\paragraph{Démo RESTO}
\begin{enumerate}
  \item Mettre le terminal en mode \textbf{RESTO} (page terminal ou endpoint \texttt{terminal\_mode\_set.php}).
  \item Passer la carte RFID : l’API vérifie carte+solde puis débite le ticket.
  \item Résultat attendu : succès + transaction enregistrée, LED verte.
\end{enumerate}

\paragraph{Démo BUVETTE}
\begin{enumerate}
  \item Mettre le terminal en mode \textbf{BUVETTE}.
  \item Scanner 2--3 produits (code-barres) pour constituer un panier.
  \item Cliquer \textit{Payer} : création d’une commande \texttt{PENDING}.
  \item Passer la carte RFID sur l’ESP32 : claim de la commande puis paiement.
  \item Résultat attendu : commande \texttt{PAYE} et solde mis à jour.
\end{enumerate}

\section{Architecture et organisation du code}
\subsection{Architecture (vue synthétique)}
\begin{itemize}
  \item \textbf{Web (navigateur)} : UI Admin/Étudiant/Terminal, scan webcam, Web Serial.
  \item \textbf{Serveur PHP} : pages + API JSON (PDO).
  \item \textbf{MySQL} : persistance (étudiants, comptes, cartes, produits, transferts, commandes terminal).
  \item \textbf{ESP32} : lecture RFID + appels HTTP (vérification carte / paiement / logs).
\end{itemize}

\begin{figure}[H]
  \centering
  \SmartFigure{archi.png}{0.92\linewidth}
  \caption{Architecture globale (optionnel).}
\end{figure}

\subsection{Organisation du dépôt}
\begin{lstlisting}
UniSmartPay/
  web/                 # Application PHP (pages + API)
    admin/             # UI admin
    etudiant/          # UI étudiant
    terminal/          # UI terminal (RESTO/BUVETTE)
    api/               # Endpoints JSON (paiement, scan, commandes...)
    config/            # Config app + PDO
    includes/          # Helpers (CSRF, auth, utilitaires)
  database/            # Schéma SQL (import initial)
  esp32/               # Firmware ESP32 (RFID + WiFi)
  assets/              # JS/CSS + libs scan (ZXing/Quagga)
  slides/              # Présentation finale (LaTeX)
  report/              # Rapport + figures
\end{lstlisting}

\subsection{Câblage ESP32 \texorpdfstring{$\leftrightarrow$}{<->} RC522}
\begin{figure}[H]
  \centering
  \SmartFigure{figures/caplage.jpg}{0.92\linewidth}
  \caption{Schéma de câblage ESP32 \texorpdfstring{$\leftrightarrow$}{<->} RC522 (SPI).}
\end{figure}

\section{Captures d’écran (interfaces)}
Cette section illustre rapidement les écrans principaux. Déposez les images dans \texttt{report/figures/} avec les noms indiqués ci-dessous.

\subsection{Accueil}
\begin{figure}[H]
  \centering
  \SmartFigure{figures/dashboard.png}{0.92\linewidth}
  \caption{Page d’accueil (accès + terminaux).}
\end{figure}

\subsection{Espace administrateur}
\begin{figure}[H]
  \centering
  \begin{subfigure}{0.48\textwidth}
    \centering
    \SmartFigure{figures/login_adm.png}{\linewidth}
    \caption{Connexion admin}
  \end{subfigure}
  \hfill
  \begin{subfigure}{0.48\textwidth}
    \centering
    \SmartFigure{figures/dashboard_adm.png}{\linewidth}
    \caption{Dashboard admin}
  \end{subfigure}

  \vspace{10pt}
  \begin{subfigure}{0.48\textwidth}
    \centering
    \SmartFigure{figures/ajout_etd.png}{\linewidth}
    \caption{Ajout étudiant}
  \end{subfigure}
  \hfill
  \begin{subfigure}{0.48\textwidth}
    \centering
    \SmartFigure{figures/charge.png}{\linewidth}
    \caption{Assignation carte + recharge}
  \end{subfigure}

  \vspace{10pt}
  \begin{subfigure}{0.70\textwidth}
    \centering
    \SmartFigure{figures/ajout_prd.png}{\linewidth}
    \caption{Ajout produit (code-barres)}
  \end{subfigure}

  \caption{Fonctionnalités administrateur (gestion des données).}
\end{figure}

\subsection{Espace étudiant}
\begin{figure}[H]
  \centering
  \begin{subfigure}{0.48\textwidth}
    \centering
    \SmartFigure{figures/login_etd.png}{\linewidth}
    \caption{Connexion étudiant}
  \end{subfigure}
  \hfill
  \begin{subfigure}{0.48\textwidth}
    \centering
    \SmartFigure{figures/dashboard_etd.png}{\linewidth}
    \caption{Dashboard (solde + transactions)}
  \end{subfigure}

  \vspace{10pt}
  \begin{subfigure}{0.92\textwidth}
    \centering
    \SmartFigure{figures/donneés de carte.png}{0.78\linewidth}
    \caption{Détails carte RFID}
  \end{subfigure}
  \caption{Fonctionnalités étudiant.}
\end{figure}

\subsection{Terminal}
\begin{figure}[H]
  \centering
  \begin{subfigure}{0.48\textwidth}
    \centering
    \SmartFigure{figures/mode_resto.png}{\linewidth}
    \caption{Terminal RESTO}
  \end{subfigure}
  \hfill
  \begin{subfigure}{0.48\textwidth}
    \centering
    \SmartFigure{figures/mode_buvette.png}{\linewidth}
    \caption{Terminal BUVETTE (scan + panier)}
  \end{subfigure}

  \vspace{10pt}
  \begin{subfigure}{0.48\textwidth}
    \centering
    \SmartFigure{figures/voir_resltat.png}{\linewidth}
    \caption{Résultat paiement RESTO}
  \end{subfigure}
  \hfill
  \begin{subfigure}{0.48\textwidth}
    \centering
    \SmartFigure{figures/voir_resltat.png}{\linewidth}
    \caption{Web Serial Monitor}
  \end{subfigure}
  \caption{Interfaces terminal et monitoring.}
\end{figure}

\section{Guide utilisateur (détaillé)}
\subsection{Administrateur (préparation des données)}
\begin{enumerate}
  \item Se connecter : \url{http://localhost/UniSmartPay/web/admin/login.php}.
  \item Créer un étudiant (ou vérifier l’étudiant démo).
  \item Assigner une carte RFID (UID) à l’étudiant.
  \begin{itemize}
    \item Récupérer l’UID depuis le Serial Monitor (ou Web Serial) côté terminal.
    \item Coller l’UID dans la page d’assignation et activer la carte.
  \end{itemize}
  \item Recharger le solde du compte (pour autoriser le débit RESTO/BUVETTE).
  \item Ajouter des produits (mode BUVETTE) : nom, prix, et \textbf{code-barres} valide.
  \item Vérifier l’historique : après paiement, contrôler la référence et les soldes avant/après.
\end{enumerate}

\subsection{Étudiant (utilisation)}
\begin{enumerate}
  \item Se connecter : \url{http://localhost/UniSmartPay/web/etudiant/login.php}.
  \item Consulter le solde, la carte RFID, et l’historique des transactions.
  \item Paiement RESTO : passer la carte sur le terminal pour débiter le ticket fixe.
  \item Paiement BUVETTE : scanner les produits pour constituer le panier, puis payer par RFID.
\end{enumerate}

\subsection{Terminal (RESTO / BUVETTE) — pas à pas}
\paragraph{RESTO (ticket fixe)}
\begin{enumerate}
  \item Mettre le terminal en mode \textbf{RESTO}.
  \item Passer la carte RFID : lecture UID, vérification côté serveur, puis débit.
  \item Résultat attendu : statut succès + écriture d’une transaction, feedback LED verte.
\end{enumerate}

\paragraph{BUVETTE (panier)}
\begin{enumerate}
  \item Mettre le terminal en mode \textbf{BUVETTE}.
  \item Scanner les produits (code-barres) pour remplir le panier (quantité/prix).
  \item Cliquer \textit{Payer} : une commande \texttt{PENDING} est créée.
  \item Passer la carte RFID sur l’ESP32 : claim de la commande puis paiement.
  \item Résultat attendu : commande \texttt{PAYE} et solde mis à jour.
\end{enumerate}

\paragraph{Mode dynamique (piloté depuis le site)}
Le mode du terminal (RESTO/BUVETTE) peut être synchronisé par la page web (mise à jour BD), puis lu par l’ESP32 au moment du paiement (endpoints \texttt{terminal\_mode\_get.php} / \texttt{terminal\_mode\_set.php}).

\section{Dépannage (problèmes fréquents)}
\begin{itemize}
  \item \textbf{Erreur MySQL} : vérifier que MySQL est démarré et que \texttt{DB\_PORT=3307} correspond à votre XAMPP.
  \item \textbf{Connexion DB instable} : préférer \texttt{DB\_HOST=127.0.0.1} (évite certains problèmes de socket/IPv6).
  \item \textbf{Caméra bloquée} : la webcam peut être bloquée en HTTP non sécurisé. Utiliser \url{http://localhost} ou activer HTTPS.
  \item \textbf{Carte introuvable} : assigner l’UID à un étudiant via l’admin et vérifier le statut \texttt{active}.
  \item \textbf{LED toujours allumée} : vérifier le câblage (GPIO→LED→GND) ou activer la logique \textit{active-low} dans le firmware.
  \item \textbf{Timeout ESP32} : vérifier l’IP serveur (pas \texttt{localhost}), le WiFi, et l’accès Apache depuis le réseau.
  \item \textbf{UID non reconnu} : vérifier le format (majuscules/minuscules, suppression des espaces) et l’enregistrement en base.
\end{itemize}

\end{document}
